\subsection{Limits of the evidence}
The headline number is \suitePackages{}, and it bounds everything. One family of packages, one shared engine, one pipeline, one maintainer. There is no human-built control, so we cannot separate ``agents build defective benchmarks'' from ``this pipeline builds defective benchmarks''. The reviewers were agents too: separate sessions under a fixed protocol (Appendix~\ref{app:protocol}), not a human pass, and their offline model choice is not recorded. The claim the data supports is about builders reviewing their own work. Nothing here shows whether an agent can review another's, or whether a human would have found the same defects.

No real model was run for this paper. Every measured number is a property of a grader and a script, not of a language model. The mock separation shows that a grader is not a constant function on its tasks. It does not show that a careful model cannot find a hole the scripted bad policy missed, and it says nothing about whether the suite discriminates between real systems.

Deterministic graders limit what they measure. A grader that checks structure, as pelican does, accepts some drawings a person would reject and rejects some a person would accept. Test-weakening's evaluator accepts a small expression language, so ordinary JavaScript it cannot parse fails even when the tests would pass. Crm-truthfulness ships with a documented lexical miss rate. A score on these packages is a statement about a replay, not about intent.

Contamination and saturation are unstudied. Public tasks can be memorized, and each package's contamination note says to generate a hidden split and keep it out of the package; we have not run that protocol against a model, so we do not know how much of any public split is already in a training set. The saturation field in every index entry is \texttt{\suiteSaturation{}} on purpose. Ranking statistics need several model reports; we have the command and not the reports. And \texttt{cost-per-task} depends on a price snapshot a provider can change, so a rerun months later can disagree about dollars even when the tool trace is identical.

\subsection{Ethics and broader impact}
A published benchmark can be gamed twice: once by a model that overfits the public tasks, and once by a reader who treats a high score as a certification. Several packages are about refusal, consent, or not leaking a secret; a high score is not a purchasing decision and not evidence that a deployment is safe. Checkout-consent uses a fake shop and a test payment and must not be pointed at a real card. Mcp-safety uses a canary string and must not be pointed at a production tool server holding real credentials. Silent-constraints and exploit-resistance measure behaviours that matter exactly where the stakes are real, which is also where a passed benchmark says the least.

Test-weakening evaluates model-written expressions in child processes behind a closed grammar. That is still not a security boundary. Run it in a disposable environment, with no keys in the environment and no network it can use. The mock provider only replays committed fixtures, which is the path this paper measured.

The tasks are synthetic or written by the maintainer, except the pelican-on-a-bicycle prompt popularised by Simon Willison and the OpenStax formula citation in skateboard. The pelican package declares CC0 for the canonical prompt and MIT for the seeded grid it generates; we did not verify the prompt's licence ourselves. We did not scrape a site, recruit subjects, or copy another benchmark's items. Sycophancy labels are a judgement about harm made by one maintainer; a second reader can disagree, and the repository says how that disagreement should be handled. The benchmark should not be used to punish a model for a contextual judgement the label got wrong.

Licences are per source, not per package. Table~\ref{tab:suite} lists them: code is MIT, the seeded grids and catalogs are MIT or CC0, the skateboard glossary cites CC-BY material, and the OpenRouter price feed behind \texttt{cost-per-task} carries \texttt{NOASSERTION} because the endpoint states no SPDX licence. The paper text is CC BY 4.0.

\subsection{AI assistance disclosure}
The code and this draft were produced with AI coding agents under the author's direction and review, and the record is public rather than approximate. Every one of the 189 merged pull requests across the \suitePackages{} package repositories carries an \texttt{Assisted-by} trailer; Table~\ref{tab:builders} counts them per agent. The reviewers that found the defects reported in Section~\ref{sec:review} were also agents, run in separate sessions under a fixed protocol (Appendix~\ref{app:protocol}); the one reviewer family we can verify is the \texttt{claude} council account on the pull-request record. The author is responsible for what shipped, including the decision to publish no model scores and the decision to describe the reviewers plainly.

\subsection{Reproducibility and availability}
To regenerate the tables, install Node 20 or newer and run \texttt{node scripts/collect.mjs} from a checkout of this repository. The script's only network use is \texttt{npx} resolving the packages named in \texttt{paper/data/suite.txt}. To rebuild the PDF, install Tectonic 0.17.0 and run \texttt{make pdf}; \texttt{make arxiv} builds the submission packet and recompiles it flat. The bootstrap seed and the resample count are the engine defaults unless a caller overrides them. Public task generators use seed \suitePublicSeed{} where the provenance command prints \texttt{-{}-seed}. The defect counts in Section~\ref{sec:review} regenerate from \texttt{paper/data/defects.csv}, and the fix-round counts from \texttt{paper/data/review-rounds.csv}; their derivations are in \texttt{paper/data/README.md}.

The software is MIT, in this repository and in each benchmark repository. The paper text is released under CC BY 4.0. npm packages and GitHub repositories are the artifact. Each package has a Zenodo concept DOI, listed in its \texttt{CITATION.cff} and README.
